\documentclass[10pt,a4paper]{amsart}
\usepackage[margin=19mm]{geometry}
\usepackage{amsmath,amssymb,mathtools}
\usepackage[scr=rsfs,cal=euler]{mathalfa}
\usepackage{xfrac}
\usepackage[unicode,hidelinks,bookmarks=false]{hyperref}
\newcommand{\ie}{i.e.\ }
\newcommand{\cf}{cf.\ }
\newcommand{\resp}{respectively\ }
\newcommand{\Subsection}{\subsection}
\providecommand{\nobreakdash}{\nobreak\hbox{-}}
\setlength{\emergencystretch}{2em}
\multlinegap0pt
\usepackage{bm}
\usepackage{bbm}
\usepackage{dsfont}
\usepackage{xspace}
\usepackage{cancel}
\usepackage{mathtools}
\usepackage{amsthm}
\makeatletter
\@ifundefined{Lemma}{\newtheorem{Lemma}{Lemma}}{}
\makeatother
\def\E{\mathbb{E}}
\def\Var{{\rm Var}}
\def\bth{{\boldsymbol{\theta}}}
\newcommand{\Pro}{\mathbb{P}}
\title[Optimal Differentially Private Ranking from Pairwise Compari]{Optimal Differentially Private Ranking from Pairwise Comparisons}
\author{T. Tony Cai, Abhinav Chakraborty,, Yichen Wang}
\date{Source: arXiv:2507.09388v1 (CC BY 4.0)}
\begin{document}
\maketitle

\noindent\emph{Source and licence.} Optimal Differentially Private Ranking from Pairwise Comparisons. Version arXiv:2507.09388v1 (CC BY 4.0), distributed under the Creative Commons Attribution 4.0 International licence, as recorded in the arXiv licence metadata for this exact version and shown on the abstract page https://arxiv.org/abs/2507.09388v1. Full rights notice: licenses/arxiv-2507.09388-CC-BY-4.0.txt.\par\smallskip

\noindent\textit{Editorial scope.} Counted unit: the Lemma whose source label is \texttt{lem:grad} in the article (source0.tex lines 1882--1900, source label \texttt{lem:grad}), delivered together with the article's own complete original proof of it (source0.tex lines 1902--1977, one \texttt{proof} environment of 76 lines). No mathematics is written, repaired or re-derived: statement and proof are the article's own text line for line. Required context retained uncounted: lem:degree (lines 1803--1825). Every definition, display and label of those lines is retained unchanged. Omitted from this package and not redistributed: 1--1802, 1826--1881, 1901--1901, 1978--2493. Nothing inside a retained excerpt is omitted or altered: every retained body is delivered exactly as the article's lines are written, so no omission receipt of any kind accompanies this package. Citations: the retained excerpts carry no citation command at all, so no bibliography entry of the article is transcribed and no reference is redistributed. Reference adaptation: none was needed and none was made; every reference inside the delivered unit resolves inside it, so no unresolved reference or citation is produced. Printed slips kept literal: no printed word, symbol, hypothesis or number of the counted statement or of its proof was altered. Layout and class: the article's own class is \texttt{article} with packages including amsmath, amssymb, amsthm, graphicx, xcolor, mathrsfs, setspace, bbm, bm, url, hyperref, natbib, varwidth, algorithm; the wrapper is the lane's pinned \texttt{amsart} editorial wrapper at 10pt on A4, which restates the theorem-like environments the retained text needs and adds the article's own macro definitions that the retained text needs (\texttt{\textbackslash E}, \texttt{\textbackslash Pro}, \texttt{\textbackslash Var}, \texttt{\textbackslash bth}), with extra wrapper packages (bm, bbm, dsfont, xspace, cancel, mathtools). The delivered numbering is the unit's own. Assets: no retained span contains a figure environment, a graphics-inclusion command, a PDF-page inclusion, a table or a TikZ picture; no image file is copied into this package, the package declares no asset dependency and no image file is redistributed with it. Licence: version arXiv:2507.09388v1 of this article is distributed under the Creative Commons Attribution 4.0 International licence, as recorded in the arXiv licence metadata for this exact version and shown on the abstract page https://arxiv.org/abs/2507.09388v1. A full rights notice is delivered at licenses/arxiv-2507.09388-CC-BY-4.0.txt. Characters: every retained excerpt is pure ASCII, so the delivered engine, XeLaTeX with its default Latin Modern text font, reports no missing character.
\begin{Lemma}[degree concentration under \textit{with}-replacement sampling]
	\label{lem:degree}
	Let
	\[
	d_i \;=\;\sum_{j\ne i}M_{ij}, 
	\qquad
	S \;=\; \E[d_i]= \frac{2mL}{n},
	\qquad
	d_{\min}=\min_{i}d_i,\;d_{\max}=\max_{i}d_i .
	\]
	Assume
	\[
	mL \;\ge\; c\,n\log n
	\quad\text{with a numerical constant }c\ge 60 .
	\]
	Then
	\[
	\Pro\!\Bigl(
	\tfrac12 S \;\le\; d_{\min}\;\le\; d_{\max}\;\le\; \tfrac32 S
	\Bigr)
	\;\ge\; 1-O\!\bigl(n^{-8}\bigr).
	\]
\end{Lemma}

\begin{Lemma}[gradient concentration at $\bth^{*}$]
	\label{lem:grad}
	Assume \textnormal{(A1)} and suppose the degree event
	$\mathcal A_{0}$ of Lemma~\ref{lem:degree} holds.  Set
	\[
	S=\frac{2mL}{n},
	\qquad
	C_{*}:=\sqrt{33}\approx 5.75 .
	\]
	If $mL\ge c\,n\log n$ for a sufficiently large absolute constant\/
	$c$, then
	\[
	\Pro\!\Bigl(
	\bigl\|\nabla\mathcal L(\bth^{*})\bigr\|_{2}
	\;>\; C_{*}\,\kappa_{1}\sqrt{nS\log n}
	\Bigr)
	\;\le\; O\!\bigl(n^{-10}\bigr).
	\]
\end{Lemma}

\begin{proof}
	\textbf{Step 0: score as a sum of i.i.d.\ vectors.}
	Index the $mL$ comparisons by $\ell=1,\dots ,mL$.  Let
	\((I_\ell,J_\ell)\) be the unordered pair drawn at step~$\ell$ and let
	\(W_\ell\in\{0,1\}\) indicate whether the smaller-indexed item wins.
	The triples
	\(\bigl(I_\ell,J_\ell,W_\ell\bigr)\) are i.i.d., and
	\(\E[W_\ell\mid I_\ell,J_\ell]=F_{I_\ell J_\ell}\).
	
	Define
	\[
	\mathbf Z_\ell
	\;=\;
	\bigl(F_{I_\ell J_\ell}-W_\ell\bigr)\,
	\frac{F'_{I_\ell J_\ell}}{F_{I_\ell J_\ell}\bigl(1-F_{I_\ell J_\ell}\bigr)}
	\bigl(\mathbf e_{I_\ell}-\mathbf e_{J_\ell}\bigr),
	\qquad
	\ell=1,\dots,mL .
	\]
	Then \(\E\mathbf Z_\ell=\mathbf 0\) and
	\(\|\mathbf Z_\ell\|_{\infty}\le\kappa_{1}\) by (A1).  Moreover
	\(
	\nabla\mathcal L(\bth^{*})=\sum_{\ell=1}^{mL}\mathbf Z_\ell.
	\)
	
	\medskip
	\noindent
	\textbf{Step 1: sub-Gaussian parameter for one coordinate.}
	Fix \(i\in[n]\).  The $i$-th coordinate of \(\mathbf Z_\ell\) is non-zero
	iff \(i\) participates in comparison~$\ell$, an event of probability
	\(2/n\).  Therefore
	\[
	\sigma_i^{2}
	:=\sum_{\ell=1}^{mL}
	\Var\!\bigl([\mathbf Z_\ell]_{i}\bigr)
	\;\le\;
	(mL)\frac{2}{n}\,\kappa_{1}^{2}
	=\kappa_{1}^{2}S .
	\]
	Thus \([\nabla\mathcal L(\bth^{*})]_{i}\) is sub-Gaussian with
	proxy variance \(\kappa_{1}^{2}S\).
	
	\medskip
	\noindent
	\textbf{Step 2: tail bound for one coordinate.}
	For any \(t>0\),
	\[
	\Pro\bigl(|g_i|>t\bigr)
	\;\le\;
	2\exp\!\Bigl(
	-\frac{t^{2}}{2\kappa_{1}^{2}S}
	\Bigr),
	\qquad
	g_i:=\bigl[\nabla\mathcal L(\bth^{*})\bigr]_i .
	\]
	Choose \(t=C_{*}\kappa_{1}\sqrt{S\log n}\) with
	\(C_{*}=\sqrt{33}\).
	Then
	\(\Pro\bigl(|g_i|>t\bigr)\le 2n^{-C_{*}^{2}/2}=2n^{-16.5}\).
	
	\medskip
	\noindent
	\textbf{Step 3: union bound over $n$ coordinates.}
	Summing over \(i=1,\dots,n\) gives
	\[
	\Pro\!\Bigl(
	\|\nabla\mathcal L(\bth^{*})\|_{2}
	> C_{*}\kappa_{1}\sqrt{nS\log n}
	\Bigr)
	\;\le\;
	2n^{-15.5}
	\;=\;
	O\!\bigl(n^{-10}\bigr),
	\]
	which is the stated result.
\end{proof}

\end{document}
